\section{Predictor Importance Analysis for SCF-CPS Imputation}\label{app:robustness}

To assess the relative contribution of each predictor variable to the quality of wealth imputation, we conduct two complementary analyses using the QRF model on the SCF data: a leave-one-out analysis and a progressive inclusion analysis. These analyses use six of the seven predictor variables from the main experimental setup; \texttt{own\_children\_in\_household} is excluded because it was not available in the SCF preprocessing for this robustness analysis. The results help evaluate which shared variables between the SCF and CPS are most informative for imputing net worth.

\subsection{Leave-one-out analysis}

We measure the impact of each predictor by excluding it from the QRF model and observing the resulting increase in average quantile loss relative to the full model (baseline loss: 1,661,723, computed on the original dollar scale rather than the asinh-transformed scale used in Section~5). Table~\ref{tab:leave_one_out} reports the results.

\begin{table}[h]
    \centering
    \caption{Leave-one-out predictor importance for QRF wealth imputation. Relative impact measures the percentage increase in average quantile loss when the predictor is excluded.}
    \label{tab:leave_one_out}
    \begin{tabular}{lrrr}
        \toprule
        Predictor Removed & Avg Quantile Loss & Loss Increase & Relative Impact \\
        \midrule
        interest\_dividend\_income & 3,596,108 & 1,934,385 & 116.4\% \\
        employment\_income & 2,219,377 & 557,654 & 33.6\% \\
        age & 2,189,845 & 528,122 & 31.8\% \\
        pension\_income & 1,875,120 & 213,397 & 12.8\% \\
        race & 1,728,080 & 66,358 & 4.0\% \\
        is\_female & 1,698,452 & 36,729 & 2.2\% \\
        \midrule
        Baseline (all predictors) & 1,661,723 & --- & --- \\
        \bottomrule
    \end{tabular}
\end{table}

Interest and dividend income is by far the most important predictor, with its removal more than doubling the quantile loss (116.4\% increase). Employment income and age contribute substantially as well (33.6\% and 31.8\% increases, respectively). Pension income has a moderate effect (12.8\%), while race and gender contribute relatively little to imputation accuracy (4.0\% and 2.2\%).

\subsection{Progressive inclusion analysis}

We complement the leave-one-out results by adding predictors one at a time, following the importance ordering determined by the leave-one-out analysis (Table~\ref{tab:leave_one_out}), and measuring the marginal improvement at each step. Table~\ref{tab:progressive_inclusion} reports the results.

\begin{table}[h]
    \centering
    \caption{Progressive predictor inclusion for QRF wealth imputation. Predictors are added in order of importance, and marginal improvement measures the reduction in average quantile loss from adding each predictor.}
    \label{tab:progressive_inclusion}
    \begin{tabular}{clrr}
        \toprule
        Step & Predictor Added & Avg Quantile Loss & Marginal Improvement \\
        \midrule
        1 & interest\_dividend\_income & 5,037,628 & --- \\
        2 & age & 2,752,098 & 2,285,529 \\
        3 & employment\_income & 1,932,222 & 819,876 \\
        4 & pension\_income & 1,737,197 & 195,025 \\
        5 & race & 1,704,094 & 33,103 \\
        6 & is\_female & 1,685,442 & 18,652 \\
        \bottomrule
    \end{tabular}
\end{table}

The progressive analysis confirms the leave-one-out findings. The first three predictors (interest/dividend income, age, and employment income) account for the majority of imputation accuracy, with diminishing returns from additional demographic variables. All six predictors together yield the lowest quantile loss (1,685,442), indicating that each variable contributes positively, though the marginal gains from race and gender are small. The small difference between this value and the leave-one-out baseline (1,661,723) reflects variation across cross-validation fold assignments and random seeds between the two analyses.

\subsection{Implications for the Conditional Independence Assumption}

The strong predictive power of the shared financial variables, particularly interest and dividend income, employment income, and pension income, is consistent with, though not sufficient evidence for, the plausibility of the CIA in the SCF-CPS matching context. These variables capture the main dimensions of household financial position available in both surveys. The fact that they collectively explain a large share of the variation in net worth suggests that the common variables $X$ may be sufficient to render the target variable $Y$ (wealth) approximately conditionally independent of variables $Z$ unique to the CPS, though this assumption remains fundamentally untestable as discussed in Section~2.2.
